\section{Goals of the work and how to read this report}
\label{sec:intro}

This report documents, step by step, how the obstacle-crossing data of Tests 101--110 were turned into a machine-learning
dataset, which models were trained on it, how they were tested, what they showed, and what the results mean for a
wearable device that warns the user about an obstacle through a vibration (haptic) cue. It is written so that every number
can be traced back: each method is given with its formula and its MATLAB implementation, and every table of results is
generated directly from the analysis reports in \vn{Results/ML Models/MATLAB/}.

\subsection{Research questions}
\label{sec:questions}

\begin{enumerate}[label=\textbf{Q\arabic*}]
  \item \textbf{How do people cross?} How do young adults adapt the movement of the leading and trailing limb to the height
  and width of an obstacle, where in the swing does this adaptation appear, and how much does it differ between people?
  \item \textbf{Can wearables recognise the obstacle?} Can obstacle height (three levels) and width (two levels) be
  recognised from the Movella Dot IMUs alone, for a person the model has never seen?
  \item \textbf{What would a device need?} Which sensors, how much information about the new user (calibration), how it
  behaves in a realistic session, without any camera, and in real time.
  \item \textbf{When and how should a haptic cue be given?} From which limb and body segment, at which moment, and with
  which threshold, so that a cue for the trailing limb arrives before that limb leaves the floor.
\end{enumerate}

\subsection{Main findings}

\begin{keybox}[Answers in one page]
\begin{description}[leftmargin=0.3cm,style=nextline,itemsep=3pt]
  \item[Q1 How people cross] Each foot is lifted by about the obstacle height (leading foot +10.1\,mm per cm, trailing +11.1\,mm per cm,
  measured by the foot sensor), through leading-thigh flexion and trailing heel-up, while the clearance stays constant. Height shows
  most in the trailing foot early in its swing and in the leading foot around its apex. Each person has a strong walking signature
  (one crossing identifies the person in 99\,\% of cases) and a personal obstacle margin. Nothing in the approach steps anticipates
  the height.
  \item[Q2 Recognition] Height (3 levels): generic models 56--73\,\%; with person scaling (no labels) 90--92\,\% for a new person;
  2.4\,cm error as a number; an unseen middle height is interpolated with 3.1\,cm error. Width: 54--66\,\% (chance 50\,\%),
  because people barely change how they cross a deeper obstacle.
  \item[Q3 Device] Two foot IMUs are enough (92.3\,\%); one labelled crossing gives 80\,\%, 10--15 unlabelled crossings 81--85\,\%;
  in a realistic mixed session the device reaches 90\,\% from crossing 16 on; without any camera 86\,\% of the crossings are found,
  paired and classified correctly.
  \item[Q4 Haptic cue] Cue the trailing limb from the leading limb, at the latest at the leading-foot apex (about 0.45\,s before the
  trailing toe leaves the floor); the leading landing is too late (0.2\,s). Detect the crossing from the leading thigh or foot
  (97--99\,\%); grade the height with personal thresholds (64--73\,\% with one signal, 71--76\,\% with the multi-feature real-time
  model); the trailing limb's own signals come too late for a cue.
\end{description}
\end{keybox}

\subsection{Overall approach}

\begin{enumerate}[leftmargin=*]
  \item \textbf{Dataset} (Section~\ref{sec:data}). One row per obstacle crossing (or level walk) with both limbs, built by
  \vn{Analysis/ML/build_crossing_master.m} from the outputs of the processing pipeline (steps 1--9). Wearable (Dot) values
  are the only model inputs; camera values are kept as reference.
  \item \textbf{Evaluation rule} (Section~\ref{sec:eval}). Every accuracy is measured on a participant left out of training
  (leave-one-participant-out). Every fitted step (imputation, scaling, feature selection, model) uses the training
  participants only.
  \item \textbf{Models} (Section~\ref{sec:models}). Logistic regression, random forests, support vector machines, boosted
  trees, ensembles, a 1D convolutional network, regression models and a physics-informed offset model.
  \item \textbf{Analyses} (Section~\ref{sec:analyses}). Nine MATLAB scripts, each answering part of Q1--Q4
  (Table~\ref{tab:analysis-map}).
  \item \textbf{Haptic design} (Sections~\ref{sec:res-haptic} and \ref{sec:discussion}). Timing of the crossing events
  and single-signal thresholds that a simple cueing rule could use.
\end{enumerate}

\subsection{How to read the numbers}

\begin{itemize}
  \item \textbf{Balanced accuracy} is the mean of the per-class recalls; it is the main metric. Chance is
  33.3\,\% for height (three classes), 50\,\% for width and 16.7\,\% for the six obstacles.
  \item \textbf{95\,\% CI} = participant bootstrap: the ten participants are resampled with replacement 2000 times.
  \item \textbf{H1 / H2 / H3} = obstacle heights of 10 / 20 / 30\,\% of the participant's leg length;
  \textbf{W1 / W2} = obstacle widths (depth along the walkway) of 5 / 15\,cm.
  \item \textbf{Leading / trailing limb} = the first / second limb to cross the obstacle.
  \item \textbf{Generic model} = the features are only referenced to the person's normal (level) walking;
  \textbf{person scaling} = each feature is additionally scaled with the person's own crossings (median and
  inter-quartile range), without using any obstacle label (Section~\ref{sec:norm}).
\end{itemize}
